\section{An Exploration of Softmax and Its Variants}
\begin{algorithm}[H]
    \caption{Single Target Loss Exploration of Softmax Scheme}
    \label{alg:velocity-add-delete}
    \textbf{Require:} Initial model $\theta_0$, reward function $R$, prompt distribution $p(c)$, number of prompts per iteration $C$, rollout size per prompt $K$, batch size $B$, total samples per iteration $D = KC$, rollout coverage coefficient $\beta$, mass-shift clip $\Delta_{\mathrm{clip}}$, reference regularization $\lambda_{\mathrm{ref}}$, EMA rate $\eta$, reward-scale constant $\epsilon_R$, numerical constant $\epsilon_{\mathrm{alg}}$, gradient accumulation steps $N_{\mathrm{acc}}$.\\
    \textbf{Initialize:} Training policy $\theta \leftarrow \theta_0$, old policy $\theta_{\mathrm{old}} \leftarrow \theta_0$, reference policy $\theta_{\mathrm{ref}} \leftarrow \theta_0$, data buffer $\mathcal{D} \leftarrow \emptyset$.
    \begin{algorithmic}[1]
        \For{\text{each iteration}}
            \For{\text{each prompt $c_j \sim p(c)$, $j = 1, \dots, C$}} \Comment{Rollout Step, Data Collection}
                \State Sample $K$ clean images $\{\mathbf{x}_0^{j,k}\}_{k=1}^{K}$ with $\theta$ conditioned on $c_j$.
                \State Evaluate rewards $R_{j,k} \leftarrow R(\mathbf{x}_0^{j,k}, c_j)$ for $k = 1, \dots, K$.
            \EndFor

            % \State Compute prompt-wise reward means
            % $\bar{R}_{c_j} \leftarrow \frac{1}{K}\sum_{k=1}^{K} R_{j,k}$ for $j = 1, \dots, C$.

            % \State Compute global reward mean
            % $\bar{R} \leftarrow \frac{1}{D}\sum_{j=1}^{C}\sum_{k=1}^{K} R_{j,k}$.

            % \State Compute global reward scale
            % $\sigma_R \leftarrow \sqrt{\frac{1}{D}\sum_{j=1}^{C}\sum_{k=1}^{K}(R_{j,k} - \bar{R})^2}.$

            \For{\text{each prompt $c_j$, $j = 1, \dots, C$}}
                \State Compute normalized weights
                \[
                    w_{j,k} \leftarrow \frac{f(R_{j,k})}{\frac{1}{G}\sum_{k=1}^{G}f(R_{j,k})},
                    \quad k = 1, \dots, K.
                \]

                \State Store
                $\mathcal{D}\leftarrow\mathcal{D}\cup
                \{(c_j,\mathbf{x}_0^{j,k},w_{j,k})\}_{k=1}^{K}$.
            \EndFor

            \For{\text{each mini batch }
            $\{(c_i,\mathbf{x}_i^0,w_i)\}_{i=1}^{B}\subset\mathcal{D}$}
                \State Sample $t_i\sim\mathcal{U}(0,1)$,
                $\boldsymbol{\epsilon}_i\sim\mathcal{N}(\mathbf{0},\mathbf{I})$;
                set $\mathbf{x}_i^{t_i}\leftarrow
                (1-t_i)\mathbf{x}_i^0+t_i\boldsymbol{\epsilon}_i$.
                \State Predict $\mathbf{x}_i^\theta$,
                $\mathbf{x}_i^{\mathrm{old}}$, and $\mathbf{x}_i^{\mathrm{ref}}$
                from $(\mathbf{x}_i^{t_i},c_i,t_i)$ with
                $\theta$, $\theta_{\mathrm{old}}$, and $\theta_{\mathrm{ref}}$.
                \State $\mathbf{x}_i^{\mathrm{target}}\leftarrow
                \mathbf{x}_i^{\mathrm{old}}+
                \beta(w_i-1)(\mathbf{x}_i^0-\mathbf{x}_i^{\mathrm{old}})$.
                \State Compute
                \[
                    \begin{aligned}
                    \mathcal{L}_{\mathrm{policy}}
                    &\leftarrow\frac{1}{B}\sum_{i=1}^{B}
                    \frac{\|\mathbf{x}_i^\theta-\mathbf{x}_i^{\mathrm{target}}\|_2^2}
                    {\operatorname{sg}[\operatorname{mean}
                    |\mathbf{x}_i^0-\mathbf{x}_i^{\mathrm{old}}|]},\\
                    \mathcal{L}_{\mathrm{ref}}
                    &\leftarrow\frac{1}{B}\sum_{i=1}^{B}
                    \|\mathbf{v}_i^\theta-\mathbf{v}_i^{\mathrm{ref}}\|_2^2.
                    \end{aligned}
                \]
                \State Backpropagate
                $(\Delta_{clip}\mathcal{L}_{\mathrm{policy}}+
                \lambda_{\mathrm{ref}}\mathcal{L}_{\mathrm{ref}})/N_{\mathrm{acc}}$.
            \EndFor
            \State Update $\theta$; set
            $\theta_{\mathrm{old}}\leftarrow
            \eta\theta_{\mathrm{old}}+(1-\eta)\theta$ and
            $\mathcal{D}\leftarrow\emptyset$.
        \EndFor
    \end{algorithmic}
    \textbf{Output:} $\mathbf{v}_\theta$
\end{algorithm}

For each prompt \(c_j\), let
\[
\bar R_{c_j}=\frac{1}{K}\sum_{k=1}^{K}R_{j,k},
\qquad
\sigma_{c_j}=\sqrt{\frac{1}{K}\sum_{k=1}^{K}\bigl(R_{j,k}-\bar R_{c_j}\bigr)^2},
\]
All weights below are normalized inside each prompt, i.e.
\[
\frac{1}{K}\sum_{k=1}^{K}w_{j,k}=1.
\]

\paragraph{Experiment A: Softmax weight variants.}

\subparagraph{Standard-deviation temperature.}
Directly substitute \(c\sigma_{c_j}\) as the temperature:
\[
f(R_{j,k})
=
\exp\left(
\frac{R_{j,k}}{c\sigma_{c_j}}
\right),
\qquad
c\in\{0.5,1.0,1.5,2.5\}.
\]

\subparagraph{Centered L1-norm temperature.}
Let
\[
\|R_{j,\cdot}-\bar R_{c_j}\|_1
=
\sum_{\ell=1}^{K}\bigl|R_{j,\ell}-\bar R_{c_j}\bigr|.
\]
Define
\[
f(R_{j,k})
=
\exp\left(
\frac{R_{j,k}}{c\,\|R_{j,\cdot}-\bar R_{c_j}\|_1}
\right),
\qquad
c\in\{0.5,1.0,1.5,2.5\}.
\]


\paragraph{Experiment B: Linear centered and uncentered weights.}

% \subparagraph{Centered linear.}
% \[
% f_{j,k}
% =
% \left[
% 1 + \frac{R_{j,k}-\bar R_{c_j}}{\sigma_{c_j}}
% \right]_+.
% \]

\subparagraph{Uncentered linear.}
\[
f_{j,k}
=
\left[
R_{j,k}
\right]_+.
\]

\subparagraph{Shifted uncentered linear.}
\[
f_{j,k}
=
\left[
1+R_{j,k}
\right]_+.
\]

where $[x]_+ = \max(x,0)$.


\paragraph{Experiment C: Soft Top-$k$.}

We consider this top-k design from a constrained optimization problem:

\begin{equation}
\max_{\{w_i\}}
\sum_{i=1}^{G} w_i R_i
\quad
\text{s.t.}
\quad
\sum_{i=1}^{G} w_i = G,
\qquad
0 \le w_i \le \frac{G}{k}.
\end{equation}


The objective assigns as much probability mass as possible to samples with higher rewards. And the upper bound $G/k$ prevents the mass focusing to fewer than $k$ samples. So the optimal solution is: samples above the boundary $\tau=R_{(k)}$ have the maximum weight $G/k$, samples below $\tau$ have no weight, and the remaining mass is equally distributed among samples equal to $\tau$.

Sort the rewards, \begin{equation} R_{(1)} \ge R_{(2)} \ge \cdots \ge R_{(G)}, \end{equation} and define the boundary reward as \begin{equation} \tau = R_{(k)}. \end{equation}

\begin{equation}
w_i=
\begin{cases}
\frac{G}{k},
& R_i>\tau,\\[4pt]
\displaystyle
\frac{G}{k}
\frac{
k-\sum_j \mathbbm{1}[R_j>\tau]
}{
\sum_j \mathbbm{1}[R_j=\tau]
},
& R_i=\tau,\\[10pt]
0,
& R_i<\tau.
\end{cases}
\end{equation}


This promises
\begin{equation}
\sum_{i=1}^{G} w_i = G.
\end{equation}

For continuous rewards with distinct values, it reduces to standard Top-$k$ weighting.


